\documentclass[11pt, a4paper]{article}
\usepackage[margin=1in]{geometry}
\usepackage[T1]{fontenc}
\usepackage{lmodern}
\usepackage{microtype}
\usepackage{booktabs}
\usepackage{tabularx}
\usepackage{amsmath, amssymb}
\usepackage{enumitem}
\usepackage{xcolor}
\usepackage{hyperref}
\usepackage{float}
\usepackage{caption}
\usepackage{subcaption}
\usepackage{multirow}
\usepackage{graphicx}

\graphicspath{{figures/}}

\definecolor{accentblue}{HTML}{2563EB}
\definecolor{muted}{HTML}{64748B}
\definecolor{neg}{HTML}{DC2626}
\definecolor{pos}{HTML}{16A34A}

\hypersetup{
  colorlinks=true,
  linkcolor=accentblue,
  urlcolor=accentblue,
  citecolor=accentblue,
}

\setlength{\parindent}{0pt}
\setlength{\parskip}{6pt}

\usepackage{titlesec}
\titleformat{\section}{\large\bfseries\sffamily}{\thesection.}{0.5em}{}[\vspace{-4pt}\rule{\textwidth}{0.4pt}]
\titleformat{\subsection}{\normalsize\bfseries\sffamily}{\thesubsection}{0.5em}{}
\titleformat{\subsubsection}{\normalsize\itshape}{\thesubsubsection}{0.5em}{}

\title{%
  \sffamily\bfseries
  Annexure~2: Multi-Shot Reconsideration Experiment\\[4pt]
  \Large Effect of Conversational Follow-Up on\\
  VLM-Based Embryo Stage Classification
}
\author{%
  Gently Project, Shroff Lab, Janelia Research Campus
}
\date{March 2026}

\begin{document}
\maketitle

\begin{abstract}
We test whether multi-turn conversational reconsideration improves VLM-based
developmental stage classification.  After an initial single-shot classification,
a follow-up prompt asks the model to re-examine the image and confirm or revise
its answer, while retaining the full conversation context (reference images,
initial reasoning, and the current image).

The result is unambiguously negative.  Multi-shot reconsideration reduces exact
accuracy by 32--36 percentage points across two prompt variants (minimal: 48.5\%
$\to$ 12.7\%; descriptive: 48.0\% $\to$ 15.5\%).  The model systematically
interprets the reconsideration prompt as implicit negative feedback, changing
correct answers to incorrect ones.  Early-stage classification is most severely
affected (98\% $\to$ 5--6\%), but no stage benefits reliably from reconsideration.
\end{abstract}

% ═══════════════════════════════════════════════════════════════
\section{Motivation}
% ═══════════════════════════════════════════════════════════════

Single-shot VLM classification of \emph{C.~elegans} developmental stages
achieves $\sim$48\% exact accuracy with both minimal and descriptive prompt
variants.  Accuracy varies widely by stage: early (98\%), bean (46--75\%), 2-fold
(46--82\%), but pretzel (29--33\%) and 1.5-fold (12--18\%) remain challenging.

A natural hypothesis is that the model might improve if given a chance to
reconsider---analogous to ``think again'' or chain-of-verification approaches in
LLM reasoning.  By continuing the conversation after the initial prediction, the
model retains access to:

\begin{itemize}[nosep]
  \item All reference stage images (cached in the first user message)
  \item Its own initial classification and reasoning
  \item The current embryo image
  \item Temporal context from previous timepoints
\end{itemize}

The follow-up prompt is intentionally neutral: it asks the model to ``carefully
re-examine the image'' and states that ``it's OK to change your answer.''  No
tools are provided on the reconsideration turn (to avoid unnecessary token cost).

% ═══════════════════════════════════════════════════════════════
\section{Experimental Setup}
% ═══════════════════════════════════════════════════════════════

\subsection{Variants}

Four variants were compared, all using \texttt{claude-sonnet-4-5} at
temperature~0.0:

\begin{center}
\small
\begin{tabular}{@{}llcl@{}}
\toprule
\textbf{Variant} & \textbf{Prompt style} & \textbf{Turns} & \textbf{Tools} \\
\midrule
Minimal (baseline)         & Stage names only         & 1 & None \\
Descriptive (baseline)     & Projection-grounded desc. & 1 & None \\
Minimal multishot          & Stage names only         & 2 & None \\
Descriptive multishot      & Projection-grounded desc. & 2 & None \\
\bottomrule
\end{tabular}
\end{center}

\subsection{Reconsideration prompt}

After the initial classification, the following follow-up message is appended to
the conversation:

\begin{quote}\small\ttfamily
You classified this as \{stage\} with \{confidence\} confidence.
Before I accept this, please carefully re-examine the image.

Look specifically at:\\
1. The key morphological features that distinguish this stage from adjacent stages\\
2. Whether your observation matches the reference examples above\\
3. All available views --- not just the primary one

After re-examining, provide your revised (or confirmed) classification
in the same JSON format. It's OK to change your answer.
\end{quote}

The prompt is designed to be neutral---it does not suggest any particular
direction of change.  The model's initial response is included as an assistant
message so the full conversation context is available.

\subsection{Dataset}

The same offline testset as the prompt ablation experiment:
4~embryos, 769~timepoints total, session \texttt{59799c78}.
Ground truth annotations by a biologist (transition-timepoint format).

% ═══════════════════════════════════════════════════════════════
\section{Results}
% ═══════════════════════════════════════════════════════════════

\subsection{Overall accuracy}

\begin{center}
\begin{tabular}{@{}lrrrrr@{}}
\toprule
\textbf{Variant} & \textbf{Exact} & \textbf{$\Delta$} & \textbf{Adjacent} & \textbf{Conf.} & \textbf{ECE} \\
\midrule
Minimal (1-shot)      & 48.5\% & ---           & 65.4\% & 0.84 & --- \\
Minimal multishot     & 12.7\% & \textcolor{neg}{$-$35.8pp} & 34.7\% & 0.84 & 0.718 \\
\addlinespace
Descriptive (1-shot)  & 48.0\% & ---           & 65.1\% & 0.82 & --- \\
Descriptive multishot & 15.5\% & \textcolor{neg}{$-$32.5pp} & 47.1\% & 0.83 & 0.691 \\
\bottomrule
\end{tabular}
\end{center}

Multi-shot reconsideration reduces exact accuracy by 32--36 percentage points.
Adjacent accuracy drops by 18--31 percentage points.  Mean confidence remains
high ($\sim$0.83--0.84), indicating the model is confidently wrong after
reconsideration.

\subsection{Per-stage accuracy}

\begin{center}
\begin{tabular}{@{}lrrrrrr@{}}
\toprule
& \multicolumn{3}{c}{\textbf{Minimal}} & \multicolumn{3}{c}{\textbf{Descriptive}} \\
\cmidrule(lr){2-4} \cmidrule(l){5-7}
\textbf{Stage} & \textbf{1-shot} & \textbf{Multi} & \textbf{$\Delta$} & \textbf{1-shot} & \textbf{Multi} & \textbf{$\Delta$} \\
\midrule
early    & 98\% &  6\% & \textcolor{neg}{$-$92pp} & 98\% &  5\% & \textcolor{neg}{$-$93pp} \\
bean     & 75\% &  0\% & \textcolor{neg}{$-$75pp} & 46\% &  0\% & \textcolor{neg}{$-$46pp} \\
comma    & 22\% &  4\% & \textcolor{neg}{$-$18pp} & 56\% & 48\% & \textcolor{neg}{$-$8pp}  \\
1.5-fold & 12\% &  0\% & \textcolor{neg}{$-$12pp} & 18\% &  4\% & \textcolor{neg}{$-$14pp} \\
2-fold   & 82\% & 22\% & \textcolor{neg}{$-$60pp} & 46\% & 61\% & \textcolor{pos}{$+$15pp} \\
pretzel  & 29\% & 16\% & \textcolor{neg}{$-$13pp} & 33\% & 11\% & \textcolor{neg}{$-$22pp} \\
\midrule
\textit{N} & \multicolumn{3}{c}{769} & \multicolumn{3}{c}{769} \\
\bottomrule
\end{tabular}
\end{center}

Only one cell shows improvement: descriptive multishot on 2-fold
($+$15\,pp).  Every other stage--variant combination is worse, often
catastrophically so.

% ═══════════════════════════════════════════════════════════════
\section{Analysis}
% ═══════════════════════════════════════════════════════════════

\subsection{The ``implicit negative feedback'' effect}

The dominant failure mode is clear: the model interprets ``please carefully
re-examine'' as a signal that its first answer was wrong, even though the prompt
explicitly states that confirming the original answer is acceptable.

This manifests as:

\begin{enumerate}[nosep]
  \item \textbf{Early $\to$ later stage shift.}  The model changes correct
    ``early'' predictions to bean, comma, 1.5-fold, or beyond.  In the minimal
    multishot confusion matrix, only 9 of 157 early-stage timepoints are
    correctly classified (vs.\ 154 in single-shot).  The model distributes
    predictions across all later stages.

  \item \textbf{Bean stage elimination.}  Bean achieves 0\% accuracy in both
    multishot variants---the model always ``upgrades'' bean to a later stage on
    reconsideration.

  \item \textbf{Confabulated morphological features.}  The model's reconsideration
    reasoning contains hallucinated observations: ``clear elongation and a
    bean-like shape'' or ``asymmetric curvature characteristic of the comma
    stage'' for images that are unambiguously early-stage (round, uniform
    blastomeres with no morphological differentiation).
\end{enumerate}

\subsection{Directional bias}

The prediction shift is almost exclusively toward later (more morphologically
complex) stages.  When asked to ``look more carefully,'' the model appears to
search for features that would justify a more specific classification, finding
pattern-matching artifacts that do not correspond to real developmental features.

\begin{center}
\small
\begin{tabular}{@{}lrrr@{}}
\toprule
\textbf{Shift direction} & \textbf{Minimal} & \textbf{Descriptive} \\
\midrule
Earlier $\to$ later (net wrong)     & dominant & dominant \\
Later $\to$ earlier (net wrong)     & rare     & rare \\
Same answer confirmed               & minority & minority \\
\bottomrule
\end{tabular}
\end{center}

This asymmetry suggests the model has a prior that ``looking more carefully''
should reveal additional structure, biasing it toward later developmental stages
that have more distinctive morphology.

\subsection{Comparison with verification subagents}

The production perception engine includes a verification mechanism that spawns
parallel subagents for binary stage comparisons (e.g., ``Is this 2-fold or
pretzel?'').  That approach was also found to hurt accuracy in the baseline
experiment (23\% with tools vs.\ 35\% without).  The multi-shot result is
consistent: \textbf{asking the same model to second-guess itself---whether via
tools, subagents, or follow-up prompts---systematically degrades performance on
this task.}

% ═══════════════════════════════════════════════════════════════
\section{Conclusion}
% ═══════════════════════════════════════════════════════════════

Multi-turn reconsideration with a generic ``are you sure?''\ prompt does not
improve VLM-based embryo classification.  The model treats conversational
follow-up as implicit negative feedback, changing correct answers with high
confidence.

This finding has implications beyond microscopy:

\begin{itemize}[nosep]
  \item \textbf{VLM self-consistency is fragile for visual classification.}
    Unlike reasoning tasks where chain-of-thought can improve accuracy,
    perceptual classification appears to degrade with deliberation pressure.

  \item \textbf{``Are you sure?'' is a biased prompt.}  Even when phrased
    neutrally, the act of asking for reconsideration carries implicit
    information that the first answer may be wrong.

  \item \textbf{First impressions may be best.}  For morphological
    classification tasks, the model's initial gestalt recognition may be more
    reliable than its ability to reason analytically about specific features
    on demand.
\end{itemize}

Future work could explore \textbf{independent multi-sample} approaches
(multiple separate classifications with majority vote) rather than
\textbf{sequential reconsideration}, which avoids the implicit negative feedback
problem.

% ═══════════════════════════════════════════════════════════════
\appendix
\section{Confusion Matrices}
% ═══════════════════════════════════════════════════════════════

\subsection*{A.1 Minimal Multishot}

\begin{center}
\small
\begin{tabular}{@{}l*{8}{r}@{}}
\toprule
& \multicolumn{8}{c}{\textbf{Predicted}} \\
\cmidrule(l){2-9}
\textbf{True} & early & bean & comma & 1.5f & 2f & pretzel & hatching & hatched \\
\midrule
early    & \textbf{9} & 10 & 13 & 16 & 56 & 39 & 12 & 2 \\
bean     & -- & \textbf{0} & 1 & -- & 7 & 10 & 4 & 2 \\
comma    & -- & -- & \textbf{1} & 1 & 16 & 9 & -- & -- \\
1.5-fold & -- & -- & -- & \textbf{0} & 22 & 21 & 6 & -- \\
2-fold   & -- & -- & -- & -- & \textbf{17} & 31 & 30 & 1 \\
pretzel  & 210 & 14 & 4 & -- & 36 & \textbf{71} & 68 & 30 \\
\bottomrule
\end{tabular}
\end{center}

\subsection*{A.2 Descriptive Multishot}

\begin{center}
\small
\begin{tabular}{@{}l*{8}{r}@{}}
\toprule
& \multicolumn{8}{c}{\textbf{Predicted}} \\
\cmidrule(l){2-9}
\textbf{True} & early & bean & comma & 1.5f & 2f & pretzel & hatching & hatched \\
\midrule
early    & \textbf{8} & 17 & 92 & 7 & 14 & 16 & -- & 3 \\
bean     & -- & \textbf{0} & 21 & 1 & 1 & 1 & -- & -- \\
comma    & -- & -- & \textbf{13} & 1 & 9 & 4 & -- & -- \\
1.5-fold & -- & -- & 15 & \textbf{2} & 20 & 12 & -- & -- \\
2-fold   & -- & -- & 2 & -- & \textbf{48} & 26 & 3 & -- \\
pretzel  & 234 & 2 & -- & 2 & 123 & \textbf{48} & 20 & 4 \\
\bottomrule
\end{tabular}
\end{center}

\subsection*{A.3 Baseline Minimal (1-shot, for comparison)}

\begin{center}
\small
\begin{tabular}{@{}l*{8}{r}@{}}
\toprule
& \multicolumn{8}{c}{\textbf{Predicted}} \\
\cmidrule(l){2-9}
\textbf{True} & early & bean & comma & 1.5f & 2f & pretzel & hatching & hatched \\
\midrule
early    & \textbf{154} & 3 & -- & -- & -- & -- & -- & -- \\
bean     & 5 & \textbf{18} & 1 & -- & -- & -- & -- & -- \\
comma    & -- & 15 & \textbf{6} & 2 & 4 & -- & -- & -- \\
1.5-fold & -- & 15 & 15 & \textbf{6} & 13 & -- & -- & -- \\
2-fold   & -- & -- & 7 & 7 & \textbf{65} & -- & -- & -- \\
pretzel  & 186 & -- & -- & -- & 69 & \textbf{124} & -- & 54 \\
\bottomrule
\end{tabular}
\end{center}

\subsection*{A.4 Baseline Descriptive (1-shot, for comparison)}

\begin{center}
\small
\begin{tabular}{@{}l*{8}{r}@{}}
\toprule
& \multicolumn{8}{c}{\textbf{Predicted}} \\
\cmidrule(l){2-9}
\textbf{True} & early & bean & comma & 1.5f & 2f & pretzel & hatching & hatched \\
\midrule
early    & \textbf{154} & 3 & -- & -- & -- & -- & -- & -- \\
bean     & 5 & \textbf{11} & 8 & -- & -- & -- & -- & -- \\
comma    & -- & 12 & \textbf{15} & -- & -- & -- & -- & -- \\
1.5-fold & -- & 18 & 22 & \textbf{9} & -- & -- & -- & -- \\
2-fold   & -- & -- & 10 & 14 & \textbf{36} & 19 & -- & -- \\
pretzel  & 82 & -- & -- & -- & 49 & \textbf{144} & -- & 158 \\
\bottomrule
\end{tabular}
\end{center}

\end{document}
